\section{Characterizing proper maps}
\label{section-proper}

\noindent
We include a section discussing the notion of a proper map in usual
topology. We define a continuous map of topological spaces to be
{\it proper} if it is universally closed and separated.
Although this matches well with the definition of a proper morphism
in algebraic geometry, this is different from the definition in Bourbaki.
With our definition of a proper map of topological spaces, the
proper base change theorem (Cohomology, Theorem
\ref{cohomology-theorem-proper-base-change})
holds without any further assumptions. Furthermore, given a morphism
$f : X \to Y$ of finite type schemes over $\mathbf{C}$ one
has: $f$ is proper as a morphism of schemes if and only if
the continuous map $f : X(\mathbf{C}) \to Y(\mathbf{C})$
on $\mathbf{C}$-points with the classical topology is proper.
This is explained in \cite[Exp. XII, Prop.\ 3.2(v)]{SGA1}
which also has a footnote pointing out that they take properness
in topology to be Bourbaki's notion with separatedness added on.

\medskip\noindent
We find it useful to have names for three distinct concepts:
separated, universally closed, and both of those together (i.e., properness).
For a continuous map $f : X \to Y$ of locally compact Hausdorff spaces
the word ``proper'' has long been used for the notion
``$f^{-1}($compact$) =$ compact'' and
this is equivalent to universal closedness for such nice spaces.
In fact, we will see the preimage condition formulated for clarity
using the word ``quasi-compact'' is equivalent to universal closedness
in general, if one includes the assumption of the map being closed. See
also \cite[Exercises 22-26 in Chapter II]{LangReal} but
beware that Lang uses ``proper'' as a synonym for ``universally closed'',
like Bourbaki does.

\begin{lemma}[Tube lemma]
\label{lemma-tube}
Let $X$ and $Y$ be topological spaces.
Let $A \subset X$ and $B \subset Y$ be quasi-compact subsets.
Let $A \times B \subset W \subset X \times Y$ with $W$
open in $X \times Y$. Then there exists opens $A \subset U \subset X$
and $B \subset V \subset Y$ such that $U \times V \subset W$.
\end{lemma}

\begin{proof}
For every $a \in A$ and $b \in B$ there exist opens
$U_{(a, b)}$ of $X$ and $V_{(a, b)}$ of $Y$ such that
$(a, b) \in U_{(a, b)} \times V_{(a, b)} \subset W$.
Fix $b$ and we see there exist a finite number $a_1, \ldots, a_n$
such that $A \subset U_{(a_1, b)} \cup \ldots \cup U_{(a_n, b)}$.
Hence
$$
A \times \{b\} \subset
(U_{(a_1, b)} \cup \ldots \cup U_{(a_n, b)}) \times
(V_{(a_1, b)} \cap \ldots \cap V_{(a_n, b)}) \subset W.
$$
Thus for every $b \in B$ there exists opens $U_b \subset X$ and
$V_b \subset Y$ such that $A \times \{b\} \subset U_b \times V_b \subset W$.
As above there exist a finite number $b_1, \ldots, b_m$ such
that $B \subset V_{b_1} \cup \ldots \cup V_{b_m}$.
Then we win because
$A \times B \subset
(U_{b_1} \cap \ldots \cap U_{b_m}) \times
(V_{b_1} \cup \ldots \cup V_{b_m})$.
\end{proof}

\noindent
The notation in the following definition may be slightly different
from what you are used to.

\begin{definition}
\label{definition-proper-map}
Let $f : X\to Y$ be a continuous map between topological spaces.
\begin{enumerate}
\item We say that the map $f$ is {\it closed}
if the image of every closed subset is closed.
\item We say that the map $f$ is {\it Bourbaki-proper}\footnote{This is the
terminology used in \cite{Bourbaki}. Sometimes this property may be
called ``universally closed'' in the literature.} if the map
$Z \times X\to Z \times Y$ is closed for any topological space $Z$.
\item We say that the map $f$ is {\it quasi-proper} if
the inverse image $f^{-1}(V)$ of every quasi-compact subset $V \subset Y$
is quasi-compact.
\item We say that $f$ is {\it universally closed} if
the map $f': Z \times_Y X \to Z$ is closed for any continuous map $g: Z \to Y$.
\item We say that $f$ is {\it proper} if $f$ is separated
and universally closed.
\end{enumerate}
\end{definition}

\noindent
The following lemma is useful later.

\begin{lemma}
\label{lemma-characterize-quasi-compact}
\begin{reference}
Combination of
\cite[I, p. 75, Lemme 1]{Bourbaki} and
\cite[I, p. 76, Corrolaire 1]{Bourbaki}.
\end{reference}
A topological space $X$ is quasi-compact if and only if the
projection map $Z \times X \to Z$ is closed for
any topological space $Z$.
\end{lemma}

\begin{proof}
(See also remark below.)
If $X$ is not quasi-compact, there exists an open covering
$X = \bigcup_{i \in I} U_i$ such that no finite
number of $U_i$ cover $X$.
Let $Z$ be the subset of the power set $\mathcal{P}(I)$ of $I$
consisting of $I$ and all nonempty finite subsets of $I$.
Define a topology on $Z$ with as a basis for the topology
the following sets:
\begin{enumerate}
\item All subsets of $Z\setminus\{I\}$.
\item For every finite subset $K$ of $I$ the set
$U_K := \{J\subset I \mid J \in Z, \ K\subset J \})$.
\end{enumerate}
It is left to the reader to verify this is the basis for a topology.
Consider the subset of $Z \times X$ defined by the formula
$$
M = \{(J, x) \mid J \in Z, \ x \in \bigcap\nolimits_{i \in J} U_i^c)\}
$$
If $(J, x) \not \in M$, then $x \in U_i$ for some $i \in J$.
Hence $U_{\{i\}} \times U_i \subset Z \times X$ is an open
subset containing $(J, x)$ and not intersecting $M$. Hence
$M$ is closed. The projection of $M$ to $Z$ is $Z-\{I\}$
which is not closed. Hence $Z \times X \to Z$ is not closed.

\medskip\noindent
Assume $X$ is quasi-compact. Let $Z$ be a topological space.
Let $M \subset  Z \times X$ be closed. Let $z \in Z$ be a point
which is not in $\text{pr}_1(M)$. By the Tube Lemma \ref{lemma-tube}
there exists an open $U \subset Z$ such that $U \times X$ is
contained in the complement of $M$. Hence $\text{pr}_1(M)$ is closed.
\end{proof}

\begin{remark}
\label{remark-lemma-literature}
Lemma \ref{lemma-characterize-quasi-compact} is a combination of
\cite[I, p. 75, Lemme 1]{Bourbaki} and
\cite[I, p. 76, Corollaire 1]{Bourbaki}.
\end{remark}

\begin{theorem}
\label{theorem-characterize-proper}
\begin{reference}
In \cite[I, p. 75, Theorem 1]{Bourbaki} you can find:
(2) $\Leftrightarrow$ (4).
In \cite[I, p. 77, Proposition 6]{Bourbaki} you can find:
(2) $\Rightarrow$ (1).
\end{reference}
Let $f: X\to Y$ be a continuous map between
topological spaces. The following conditions are equivalent:
\begin{enumerate}
\item The map $f$ is quasi-proper and closed.
\item The map $f$ is Bourbaki-proper.
\item The map $f$ is universally closed.
\item The map $f$ is closed and $f^{-1}(y)$ is quasi-compact for any
$y\in Y$.
\end{enumerate}
\end{theorem}

\begin{proof}
(See also the remark below.)
If the map $f$ satisfies (1), it automatically satisfies (4) because
any single point is quasi-compact.

\medskip\noindent
Assume map $f$ satisfies (4).
We  will prove it is universally closed, i.e., (3) holds.
Let $g : Z \to Y$ be a continuous map of topological spaces
and consider the diagram
$$
\xymatrix{
Z \times_Y X \ar[r]_{g'} \ar[d]_{f'} & X \ar[d]^f \\
Z \ar[r]^g & Y
}
$$
During the proof we will use that $Z \times_Y X \to Z \times X$
is a homeomorphism onto its image, i.e., that we may identify
$Z \times_Y X$ with the corresponding subset of $Z \times X$ with
the induced topology.
The image of $f' : Z \times_Y X \to Z$ is
$\Im(f') = \{z : g(z) \in f(X)\}$.
Because $f(X)$ is closed, we see that
$\Im(f')$ is a closed subspace of $Z$.
Consider a closed subset $P \subset Z \times_Y X$.
Let $z \in Z$, $z \not \in f'(P)$.
If $z \not \in \Im(f')$, then $Z \setminus \Im(f')$
is an open neighbourhood which avoids $f'(P)$.
If $z$ is in $\Im(f')$
then $(f')^{-1}\{z\} = \{z\} \times f^{-1}\{g(z)\}$
and $f^{-1}\{g(z)\}$
is quasi-compact by assumption. Because $P$ is a closed
subset of $Z \times_Y X$, we have a closed $P'$ of $Z \times X$ such
that $P = P' \cap Z \times_Y X$.
Since $(f')^{-1}\{z\}$ is a subset of $P^c = P'^c \cup (Z \times_Y X)^c$,
and since $(f')^{-1}\{z\}$ is disjoint from $(Z \times_Y X)^c$
we see that $(f')^{-1}\{z\}$ is contained in $P'^c$.
We may apply the Tube Lemma \ref{lemma-tube} to
$(f')^{-1}\{z\} = \{z\} \times f^{-1}\{g(z)\}
\subset (P')^c \subset Z \times X$.
This gives $V \times U$ containing
$(f')^{-1}\{z\}$ where $U$ and $V$ are open sets in $X$ and $Z$
respectively and $V \times U$ has empty intersection with $P'$.
Then the set $V \cap g^{-1}(Y-f(U^c))$ is open in $Z$ since $f$
is closed, contains $z$, and has empty intersection with the image of $P$.
Thus $f'(P)$ is closed. In other words, the map $f$ is universally closed.

\medskip\noindent
The implication (3) $\Rightarrow$ (2) is trivial.
Namely, given any topological space $Z$ consider the projection
morphism $g : Z \times Y \to Y$. Then it is easy to see
that $f'$ is the map $Z \times X \to Z \times Y$, in other
words that $(Z \times Y) \times_Y X = Z \times X$. (This identification
is a purely categorical property having nothing to do with
topological spaces per se.)

\medskip\noindent
Assume $f$ satisfies (2). We will prove it satisfies (1).
Note that $f$ is closed as $f$ can be identified with the map
$\{pt\} \times X \to \{pt\} \times Y$ which is assumed closed.
Choose any quasi-compact subset $K \subset Y$.
Let $Z$ be any topological space.
Because $Z \times X \to Z \times Y$ is closed
we see the map $Z \times f^{-1}(K) \to Z \times K$
is closed (if $T$ is closed in $Z \times f^{-1}(K)$, write
$T = Z \times f^{-1}(K) \cap T'$ for some closed
$T' \subset Z \times X$). Because $K$ is quasi-compact,
$K \times Z\to Z$ is closed by Lemma \ref{lemma-characterize-quasi-compact}.
Hence the composition $Z \times f^{-1}(K)\to Z \times K \to Z$
is closed and therefore $f^{-1}(K)$ must be quasi-compact
by Lemma \ref{lemma-characterize-quasi-compact} again.
\end{proof}

\begin{remark}
\label{remark-proof-literature}
Here are some references to the literature.
In \cite[I, p. 75, Theorem 1]{Bourbaki} you can find:
(2) $\Leftrightarrow$ (4).
In \cite[I, p. 77, Proposition 6]{Bourbaki} you can find:
(2) $\Rightarrow$ (1).
Of course, trivially we have (1) $\Rightarrow$ (4).
Thus (1), (2) and (4) are equivalent.
The equivalence of (3) and (4) is \cite[Chapter II, Exercise 25]{LangReal}.
\end{remark}

\begin{lemma}
\label{lemma-closed-map}
\begin{slogan}
A map from a compact space to a Hausdorff space is universally closed.
\end{slogan}
Let $f : X \to Y$ be a continuous map of topological spaces.
If $X$ is quasi-compact and $Y$ is Hausdorff, then $f$ is
universally closed.
\end{lemma}

\begin{proof}
Since every point of $Y$ is closed, we see from
Lemma \ref{lemma-closed-in-quasi-compact}
that the closed subset $f^{-1}(y)$ of $X$ is quasi-compact for all $y \in Y$.
Thus, by Theorem \ref{theorem-characterize-proper}
it suffices to show that $f$ is closed.
If $E \subset X$ is closed, then it is quasi-compact
(Lemma \ref{lemma-closed-in-quasi-compact}),
hence $f(E) \subset Y$ is quasi-compact
(Lemma \ref{lemma-image-quasi-compact}),
hence $f(E)$ is closed in $Y$
(Lemma \ref{lemma-quasi-compact-in-Hausdorff}).
\end{proof}

\begin{lemma}
\label{lemma-bijective-map}
Let $f : X \to Y$ be a continuous map of topological spaces.
If $f$ is bijective, $X$ is quasi-compact, and $Y$ is Hausdorff,
then $f$ is a homeomorphism.
\end{lemma}

\begin{proof}
It suffices to prove $f$ is closed, because this implies that
$f^{-1}$ is continuous. If $T \subset X$ is closed, then $T$
is quasi-compact by Lemma \ref{lemma-closed-in-quasi-compact},
hence $f(T)$ is quasi-compact by Lemma \ref{lemma-image-quasi-compact},
hence $f(T)$ is closed by Lemma \ref{lemma-quasi-compact-in-Hausdorff}.
\end{proof}













